% Level 2, Week 3 -- Lists: Collections of Things (student workbook)
\chapterweek{3}{Lists: Collections of Things}{A list keeps many values in order, and the first one is number 0.}

\section{What you will learn}
\begin{itemize}[leftmargin=1.2em]
  \item How to keep many values in one \textbf{list}.
  \item How to find an item by its \textbf{index}, starting from 0.
  \item How to add items with \texttt{append} and count them with \texttt{len}.
  \item How to visit every item with a \texttt{for} loop.
  \item How to pick colors that work for everyone.
\end{itemize}

\section{The big idea: a row of lockers}
A list is like a row of lockers. Each locker has a number painted on it, called its \textbf{index}.
The first locker is number \textbf{0}.

\begin{center}
%% alt: A list called backpack holds four items: "lantern" at index 0, "rope" at index 1, "map" at index 2 and "apple" at index 3. Asking for index 4 gives an IndexError because there is no fifth item.
\begin{tikzpicture}[font=\small,
  cell/.style={draw=black,thick,minimum width=2.2cm,minimum height=1cm,fill=kcsky!15}]
  \foreach \thing [count=\i from 0] in {lantern,rope,map,apple}{
    \node[cell] (c\i) at (\i*2.2,0) {\texttt{"\thing"}};
    \node[below=2pt of c\i] {index \textbf{\i}};
  }
  \node[left=4pt of c0,font=\small\ttfamily] {backpack =};
  \node[cell,dashed,fill=white,right=0pt of c3] (c4) {?};
  \node[below=2pt of c4,align=center] {index 4:\\\textbf{IndexError}};
\end{tikzpicture}
\end{center}

\begin{pycode}[Making and using a list]
backpack = ["lantern", "rope", "map"]
backpack.append("apple")    # adds to the END of the list
print(len(backpack))        # 4   how many items
print(backpack[0])          # lantern   the first item
print(backpack[-1])         # apple   -1 means the last item
\end{pycode}
A minus sign in the brackets means \enquote{count from the end}: \texttt{backpack[-1]} is the last item and
\texttt{backpack[-2]} is the one before it.
With 4 items, the index numbers are 0, 1, 2 and 3. Asking for \texttt{backpack[4]} crashes with
\texttt{IndexError: list index out of range}. This is called an \textbf{off-by-one} bug.

\subsection{Two ways to loop over a list}
\begin{center}
\renewcommand{\arraystretch}{1.3}
\begin{tabular}{@{}p{0.40\linewidth}p{0.56\linewidth}@{}}
\toprule
\textbf{Visit every item} & \textbf{Visit every index number} \\ \midrule
\texttt{for item in backpack:} & \texttt{for i in range(len(backpack)):} \\
\texttt{\ \ \ \ print(item)} & \texttt{\ \ \ \ print(i, backpack[i])} \\
Use it when you only need the items. & Use it when you also need the position number. \\ \bottomrule
\end{tabular}
\end{center}
Remember \texttt{random.choice} from Level 1? It picks one random item from a list:
\texttt{random.choice(backpack)} might give back \texttt{"map"}. You will use it in this week's homework.
When you need a loop counter but never use it, Python coders name it \texttt{\_}: \texttt{for \_ in range(3):} just repeats 3 times.

\subsection{Colors for everyone}
Some people see colors differently, often red and green. These eight \textbf{Okabe-Ito} colors stay different for almost everyone.
Use them, and always add a label, number or shape too, so color is never the only clue.

\begin{center}
\renewcommand{\arraystretch}{1.3}
\begin{tabular}{@{}cll@{\hspace{1.5em}}cll@{}}
\swatch{kcorange} & orange & \texttt{\#E69F00} & \swatch{kcblue} & blue & \texttt{\#0072B2} \\
\swatch{kcsky} & sky blue & \texttt{\#56B4E9} & \swatch{kcvermillion} & vermillion & \texttt{\#D55E00} \\
\swatch{kcgreen} & bluish green & \texttt{\#009E73} & \swatch{kcpurple} & reddish purple & \texttt{\#CC79A7} \\
\swatch{kcyellow} & yellow & \texttt{\#F0E442} & \swatch{kcblack} & black & \texttt{\#000000} \\
\end{tabular}
\end{center}

\begin{exercise}[Palette Stripes]
Folder: \texttt{exercises/week03/ex-palette-stripes}. Press \menu{Run}: only one stripe appears.
\begin{enumerate}[leftmargin=1.4em]
  \item TODO 1: add blue and vermillion with \texttt{palette.append(...)}. (TODO 2 is an extra: a sixth color.)
  \item TODO 3: draw a stripe for every color with \texttt{for i in range(len(colors)):}. Label each stripe \texttt{str(i)} and move \texttt{x} along.
        (A label must be text. \texttt{str()} turns a number into text: \texttt{str(3)} gives \texttt{"3"}.)
  \item TODO 4: print the first color (\texttt{palette[0]}) and the last color (\texttt{palette[-1]}).
  \item My palette has \rule{1cm}{0.4pt} colors, so the labels go from 0 to \rule{1cm}{0.4pt}.
\end{enumerate}
Done looks like: a row of 5 (or more) colored stripes, labeled 0, 1, 2, 3, 4 from left to right.
\end{exercise}

\begin{exercise}[Adventure Inventory]
Folder: \texttt{exercises/week03/ex-inventory}. Press \menu{Run}. It crashes!
\begin{enumerate}[leftmargin=1.4em]
  \item Read the error from the bottom up. Error name: \rule{3.5cm}{0.4pt}\quad In which function? \rule{3.5cm}{0.4pt}
  \item Draw the backpack \texttt{["lantern", "rope", "map", "apple"]} as boxes with their index numbers:
  \answerspace{1.4cm}
  \item TODO 3: fix \texttt{last\_item()}. TODO 1: \texttt{append} in \texttt{add\_item()}. TODO 2: loop in \texttt{has\_item()}.
  \item Check your work with \texttt{check\_my\_work.py}.
\end{enumerate}
\end{exercise}

\begin{exercise}[Top 3 High Scores]
Folder: \texttt{exercises/week03/ex-high-scores}. Two new tricks:
\texttt{sorted(scores, reverse=True)} makes a new list, biggest first, and \texttt{best[:3]} takes the first three items.
\begin{enumerate}[leftmargin=1.4em]
  \item Predict: the top 3 of \texttt{[120, 45, 300, 80]} are \rule{4cm}{0.4pt}
  \item TODO 1: sort, then take three. TODO 2: print \texttt{1st: 300}, \texttt{2nd: 120}, \texttt{3rd: 80} using the \texttt{places} list.
  \item Check your work with \texttt{check\_my\_work.py}.
\end{enumerate}
\end{exercise}

\begin{vocab}
  \vocabword{list}{Many values kept in order inside square brackets, like \texttt{["lantern", "rope"]}.}
  \vocabword{item}{One value inside a list.}
  \vocabword{index}{The position number of an item. The first item has index 0.}
  \vocabword{append}{Add an item to the end of a list: \texttt{backpack.append("map")}.}
  \vocabword{len}{Counts the items in a list: \texttt{len(backpack)}.}
  \vocabword{IndexError}{The error you get when you ask for an index that does not exist.}
  \vocabword{off-by-one}{A bug where a number is one too big or one too small, like using \texttt{len} as the last index.}
  \vocabword{accessible}{Designed so that everyone can use it, including people who see colors differently.}
\end{vocab}

\begin{reflection}
With a list of 10 items, what is the index of the last item? Why?
\answerlines{1}
Describe a list you could use in a game you would like to make:
\answerlines{2}
\end{reflection}
